\hypertarget{concurrency-mechanics-the-global-interpreter-lock}{%
\chapter{CONCURRENCY MECHANICS \& THE GLOBAL INTERPRETER
LOCK}\label{concurrency-mechanics-the-global-interpreter-lock}}

\hypertarget{the-global-interpreter-lock-gil-architecture}{%
\subsection{10.1 The Global Interpreter Lock (GIL)
Architecture}\label{the-global-interpreter-lock-gil-architecture}}

The CPython interpreter uses a global lock called the \textbf{Global
Interpreter Lock (GIL)} to ensure thread safety. * \textbf{Why CPython
requires the GIL}: CPython's internal memory management (reference
counting) and shared resources (like class and module namespaces
dictionaries) are not thread-safe. Without the GIL, concurrent
modifications by multiple threads could cause memory corruption. *
\textbf{CPython Thread Model}: CPython threads are standard OS threads
managed by the host operating system's scheduler. However, a thread must
hold the GIL to execute Python bytecode.

\begin{lstlisting}
Multi-Threading Execution under the GIL:
Thread 1: [ Acquires GIL ] --> [ Executes Bytecode ] --> [ Releases GIL ]
                                                               |
Thread 2:                       [ Waiting for GIL ] ------> [ Acquires GIL ] --> [ Executes Bytecode ]
\end{lstlisting}

\hypertarget{gil-release-and-acquisition-intervals}{%
\subsection{10.2 GIL Release and Acquisition
Intervals}\label{gil-release-and-acquisition-intervals}}

\begin{itemize}
\tightlist
\item
  \textbf{The Switch Interval}: CPython threads do not hold the GIL
  indefinitely. The interpreter forces the running thread to release the
  GIL after a set interval (configured via
  \passthrough{\lstinline!sys.setswitchinterval(interval)!}, defaulting
  to 5 milliseconds).
\item
  \textbf{Voluntary Release}: The GIL is automatically released during
  blocking operations:

  \begin{itemize}
  \tightlist
  \item
    POSIX System I/O calls (\passthrough{\lstinline!read!},
    \passthrough{\lstinline!write!}, \passthrough{\lstinline!select!}).
  \item
    Standard compression/hashing C routines.
  \item
    Custom C/C++ extensions explicitly wrapping blocks with the macros
    \passthrough{\lstinline!Py\_BEGIN\_ALLOW\_THREADS!} and
    \passthrough{\lstinline!Py\_END\_ALLOW\_THREADS!}.
  \end{itemize}
\end{itemize}

\hypertarget{threading-vs.-multiprocessing-memory-layouts}{%
\subsection{10.3 Threading vs.~Multiprocessing memory
layouts}\label{threading-vs.-multiprocessing-memory-layouts}}

\begin{itemize}
\tightlist
\item
  \textbf{Threading (\passthrough{\lstinline!threading!})}: Threads
  share a single CPython interpreter instance and the same virtual
  memory space (heap). This allows fast data sharing but restricts
  execution to a single core due to the GIL.
\item
  \textbf{Multiprocessing (\passthrough{\lstinline!multiprocessing!})}:
  Spawns separate OS processes, each with its own independent CPython
  interpreter and memory space (with its own GIL). This bypasses the GIL
  to run computations in parallel across multiple CPU cores, but
  requires Inter-Process Communication (IPC) (like pipes, queues, or
  shared memory structures) to exchange data.
\end{itemize}

\begin{center}\rule{0.5\linewidth}{0.5pt}\end{center}
